\subsection{连续函数的 Fourier 级数}
\label{sec:ch1-continuous-fourier-series}

若 $f$ 在 $x$ 的一个邻域内满足 Lipschitz 型条件, 即对某个 $\alpha>0$, 当 $|t|<\delta$ 时有 $|f(x+t)-f(x)|\le C|t|^\alpha$, 则 Dini 判别法适用于 $f$. 然而, 连续函数不一定满足这一条件, 也不一定满足我们已经见过的任何其他收敛判别法. P. du Bois--Reymond 于 1873 年给出的下述结果表明, 情形必然如此.

\begin{Theorem}
\label{thm:ch1-continuous-fourier-divergence}
存在一个连续函数, 其 Fourier 级数在某一点发散.
\end{Theorem}

Du Bois--Reymond 构造了一个具有这一性质的函数, 但这里将应用一致有界原理来证明这种函数存在; 该原理也称为 Banach--Steinhaus 定理.

\begin{Lemma}[一致有界原理]
\label{lem:ch1-uniform-boundedness}
设 $X$ 是 Banach 空间, $Y$ 是赋范向量空间, $\{T_a\}_{a\in A}$ 是从 $X$ 到 $Y$ 的一族有界线性算子. 则或者
\[
 \sup_a\|T_a\|<\infty,
\]
或者存在 $x\in X$, 使得
\[
 \sup_a\|T_ax\|_Y=\infty.
\]
\end{Lemma}

\hypertarget{chapter-01-page-019}{}
回忆 $T_a$ 的算子范数为 $\|T_a\|=\sup\{\|T_ax\|_Y:\|x\|_X\le1\}$. 这一结果的证明例如可见 Rudin \MBUnresolvedReference{citation}{[14, Chapter 5]}.

现在取 $X=C(\mathbb T)$, 赋予范数 $\|\cdot\|_\infty$, 并取 $Y=\mathbb C$. 定义 $T_N:X\to Y$ 为
\[
 T_Nf=S_Nf(0)=\int_{-1/2}^{1/2}f(t)D_N(t)\,dt.
\]
定义 Lebesgue 数
\[
 L_N=\int_{-1/2}^{1/2}|D_N(t)|\,dt;
\]
立即有 $|T_Nf|\le L_N\|f\|_\infty$. $D_N(t)$ 只有有限个零点, 因而 $\operatorname{sgn}D_N(t)$ 只有有限个跳跃间断点. 在每个间断点的一个小邻域内修改它, 可以构造连续函数 $f$, 使得 $\|f\|_\infty=1$ 且 $|T_Nf|\ge L_N-\varepsilon$. 因此 $\|T_N\|=L_N$. 于是, 如果能证明当 $N\to\infty$ 时 $L_N\to\infty$, 则由一致有界原理, 存在连续函数 $f$ 满足
\[
 \limsup_{N\to\infty}|S_Nf(0)|=\infty;
\]
也就是说, $f$ 的 Fourier 级数在 $0$ 处发散.

\begin{Lemma}
\label{lem:ch1-lebesgue-number-asymptotic}
有
\[
 L_N=\frac4{\pi^2}\log N+O(1).
\]
\end{Lemma}

\begin{Proof}
\[
\begin{aligned}
 L_N
 &=2\int_0^{1/2}\left|\frac{\sin(\pi(2N+1)t)}{\pi t}\right|dt+O(1)\\
 &=2\int_0^{N+1/2}\left|\frac{\sin(\pi t)}{\pi t}\right|dt+O(1)\\
 &=2\sum_{k=0}^{N-1}\int_k^{k+1}\left|\frac{\sin(\pi t)}{\pi t}\right|dt+O(1)\\
 &=\frac2\pi\sum_{k=0}^{N-1}\int_0^1\frac{|\sin(\pi t)|}{t+k}\,dt+O(1).
\end{aligned}
\]
\hypertarget{chapter-01-page-020}{}
继续可得
\[
\begin{aligned}
 L_N
 &=\frac2\pi\int_0^1|\sin(\pi t)|\sum_{k=1}^{N-1}\frac1{t+k}\,dt+O(1)\\
 &=\frac4{\pi^2}\log N+O(1).
\end{aligned}
\]
\end{Proof}

